\subsection{Parameter-Efficient Fine-Tuning and the Spectral Shift}
Low-Rank Adaptation (LoRA)~\citep{hu2022lora} introduced an efficient way to fine-tune pretrained models by learning low-rank additive updates while keeping the backbone frozen. Subsequent methods such as QLoRA~\citep{dettmers2023qlora}, VeRA~\citep{kopiczko2023vera}, DoRA~\citep{liu2024dora}, and FourierFT~\citep{gao2024parameter} further improved the parameter, memory, or update structure of PEFT methods. While these works primarily focus on efficient adaptation, recent literature has increasingly turned to the spectral properties of pretrained weights to guide adapter design.

Methods such as AdaLoRA~\citep{zhang2023adalora}, PiSSA~\citep{meng2024pissa}, MiLoRA~\citep{wang2025milora}, LoRA-XS~\citep{balazy2024lora}, SVFT~\citep{lingam2024svft}, and CorDA~\citep{yang2024corda} leverage Singular Value Decomposition (SVD), spectral importance, or decomposition-based initialization to allocate adaptation capacity more effectively. In parallel, RandLoRA~\citep{albert2025randlora} provides statistical guarantees for the expressiveness of random projections, motivated by the empirical success of random-basis methods such as NoLA~\citep{koohpayegani2023nola} and VeRA~\citep{kopiczko2023vera}.

Building on the empirical insight of LoRA that fine-tuning updates amplify task-relevant directions already present in pretrained weights~\citep{hu2022lora}, we pursue a different spectral trajectory. Rather than approximating expressiveness through random bases or using SVD mainly for initialization or rank allocation, UIOrthoLoRA directly parameterizes adaptation in pretrained singular coordinates. By preserving selected pretrained singular directions and introducing compact orthogonal rotations in the low-spectrum subspace, our method provides an explicit structural mechanism for balancing adaptation with preservation.

\subsection{Foundations of Generalization.}
The empirical success of low-rank adaptation is closely connected to intrinsic dimensionality. \citet{li2018measuring} introduced methods for measuring the intrinsic dimension of neural objective landscapes, and \citet{aghajanyan2020intrinsic} showed that large pretrained language models can often be fine-tuned effectively in surprisingly low-dimensional subspaces. This supports the view that successful adaptation may require only a restricted task-specific subspace rather than updates over the full parameter space.

In parallel, PAC-Bayesian analyses connect generalization to perturbation size and spectral properties of neural networks. \citet{neyshabur2017pac, neyshabur2018pacbayesian} derive bounds involving spectral norms and layer-wise perturbations, suggesting that controlling the geometry and magnitude of weight updates is central to generalization. Our approach is motivated by the intersection of these perspectives: UIOrthoLoRA restricts adaptation to structured subspaces derived from the pretrained spectrum, while its orthogonality and scaler-induced mixing analysis provide explicit control over how the update interacts with preserved pretrained singular structure.



